% Chapter 1 -- Rubric: PO1 (1.1), PO2 (problem identification), PO3 (3.2 objectives)
\chapter{Problem Identification and Formulation}
\label{ch:problem}

\section{Background and Motivation}
\label{sec:background}

A software organisation runs each of its projects through several specialised
tools: a version-control platform for source code and code review, an
issue-tracking platform for planning and sprint management, a static
code-analysis service for code quality, and a CI/CD service for builds and
deployments. Each tool exposes its own dashboard, reports its own metrics in its
own vocabulary, and answers questions only about its own slice of the work.
None of them answers the question that engineering leadership actually asks:
\emph{is this project, taken as a whole, healthy, and is it improving or
deteriorating?}

The consequence, as described by Spectrum Software \& Consulting Pvt.\ Ltd.\
(SSCL), the industry partner that originated this project, is that assessing the state of a single project means opening
four to five dashboards in turn and reconciling their readings by judgement.
Doing this across a whole portfolio is impractical, so it is done
infrequently. As a result, a project that is gradually deteriorating is often
noticed only after the damage has become visible as missed deadlines or failing
builds.

Two further gaps compound this. First, tool metrics describe what the code and
the pipeline are doing, but not how the team is doing; low morale, persistent
blockers, and loss of confidence in a plan are often the earliest signals of a
project in trouble, and they appear in no dashboard. Second, when a manager does
identify a problem and acts on it, the action is recorded nowhere, so there is
no later way to tell whether the intervention worked.

Pulse addresses these gaps. It ingests metrics from the separate tools,
normalises them into comparable health scores on a single scale, keeps every
score explainable in terms of the metrics that produced it, collects anonymous
team feedback that metrics cannot capture, and records the management actions
taken in response so that their effectiveness can be reviewed. It is designed as
a trend indicator, revealing how projects move over weeks and months, rather
than as a real-time monitor.

\section{Problem Statement}
\label{sec:problem-statement}

Engineering leadership and project managers at software organisations such as
SSCL oversee many projects at once, and each project produces its data in
several independent tools: version control, issue tracking, static code
analysis and CI/CD. Every tool reports its own metrics on its own scale, and
none of them offers a common notion of whether a project is healthy. To judge a
project's condition, a manager has to gather readings from each tool separately
and reconcile them by hand, and the tools capture nothing about how the team
itself is faring or about what management did when a problem was found.

Because this reconciliation is slow and manual, it is done infrequently and
inconsistently; projects cannot be compared on a like-for-like basis; and a
project that is gradually deteriorating is usually recognised only once the
damage has become visible. The early warning signs that only the team is aware of, such
as people losing motivation or work that keeps getting stuck, never reach any
dashboard, and because
interventions are not recorded, whether they helped is never known. What is
missing is a single, comparable and explainable view of project health that
draws on all of these tools, that captures the team's own signal without
exposing individuals, and that connects management action to measured outcome.
Closing that gap is the problem this project addresses.

\section{Complexity of the Problem}
\label{sec:complexity}

\begin{longtable}{
  >{\raggedright\arraybackslash}p{0.05\textwidth}
  >{\raggedright\arraybackslash}p{0.38\textwidth}
  >{\raggedright\arraybackslash}p{0.45\textwidth}
}
\caption{Complex engineering problem characteristics (P1--P7) and evidence from this project}
\label{tab:complexity}\\
\toprule
\textbf{Code} & \textbf{Definition} & \textbf{Evidence from this project} \\
\midrule
\endfirsthead
\toprule
\textbf{Code} & \textbf{Definition} & \textbf{Evidence from this project} \\
\midrule
\endhead
\bottomrule
\endlastfoot
P1 & \textbf{Depth of knowledge required.} Cannot be resolved without in-depth engineering knowledge at the level of one or more of K3, K4, K5, K6, or K8, which allows a fundamentals-based, first-principles analytical approach. &
K3: the weighted-normalisation scoring model, its handling of missing inputs, and its exact decomposition into per-metric contributions.
K4 and K6: asynchronous job processing with retries and idempotent scheduling; rate-limit-aware integration with third-party REST and GraphQL APIs; credential encryption, password hashing and session security.
K5: the connector abstraction that isolates provider-specific ingestion from scoring, and the snapshot-based data model that makes every score reproducible. \\
\addlinespace
P2 & \textbf{Range of conflicting requirements.} Involves wide-ranging or conflicting technical, engineering, and other issues. &
\guidance{To be completed after discussion with the team} \\
\addlinespace
P3 & \textbf{Depth of analysis required.} Has no obvious solution and requires abstract thinking and originality in analysis to formulate suitable models. &
There is no established mapping from heterogeneous tool metrics to a single comparable scale. The team had to derive, per metric, a normalisation function and weight, define how absent metrics are treated, and design the scores so that they remain exactly explainable; and had to design survey aggregation that yields team-level insight without any per-respondent record. \\
\addlinespace
P4 & \textbf{Familiarity of issues.} Involves infrequently encountered issues. &
\guidance{To be completed after discussion with the team} \\
\addlinespace
P6 & \textbf{Extent of stakeholder involvement and conflicting requirements.} Involves diverse groups of stakeholders with widely varying needs. &
Seven stakeholder groups (\Cref{tab:stakeholders}): the client's leadership, its managers, its developers as data subjects, the tenant administrator, the academic supervisor, external service providers as constraint setters, and the development team, whose needs conflict as set out under P2. \\
\addlinespace
P7 & \textbf{Interdependence.} Is a high-level problem including many component parts or sub-problems. &
Authentication and tenancy; workspace and connector configuration; metric ingestion and scheduling; scoring; survey generation, distribution and analysis; and the action log. These are dependent: scoring requires ingestion, survey generation is driven by scores and their trends, actions are recorded against both, and scheduling must space projects that share a credential to stay within rate limits. \\
\end{longtable}

\begin{longtable}{
  >{\raggedright\arraybackslash}p{0.08\textwidth}
  >{\raggedright\arraybackslash}p{0.80\textwidth}
}
\caption{Knowledge profiles referenced by P1}
\label{tab:knowledge-profile}\\
\toprule
\textbf{Code} & \textbf{Description} \\
\midrule
\endfirsthead
\toprule
\textbf{Code} & \textbf{Description} \\
\midrule
\endhead
\bottomrule
\endlastfoot
K3 & A systematic, theory-based formulation of engineering fundamentals required in the engineering discipline. \\
\addlinespace
K4 & Engineering specialist knowledge that provides theoretical frameworks and bodies of knowledge for the accepted practice areas in the engineering discipline; much is at the forefront of the discipline. \\
\addlinespace
K5 & Knowledge that supports engineering design in a practice area. \\
\addlinespace
K6 & Knowledge of engineering practice (technology) in the practice areas in the engineering discipline. \\
\addlinespace
K8 & Engagement with selected knowledge in the research literature of the discipline. \\
\end{longtable}

\section{Project Objectives and Scope}
\label{sec:objectives}

The project set out to achieve the following objectives:

\begin{enumerate}
  \item \textbf{Consolidate} the metrics of a project's version-control,
  issue-tracking, code-analysis and CI/CD tools into a single dashboard, so
  that routine status assessment needs none of the tools' own dashboards.
  \item \textbf{Normalise} those heterogeneous metrics into seven health scores
  on a common 0--100 scale, plus an overall score, so that any two projects
  can be compared directly.
  \item \textbf{Explain} every score on request: which metrics produced it,
  their measured values, their weights, and each one's contribution.
  \item \textbf{Capture} qualitative team feedback through anonymous surveys,
  with no stored link between a respondent and their answers.
  \item \textbf{Record} the management actions taken in response to identified
  problems, and prompt for a review of their effectiveness after a set
  interval.
  \item \textbf{Accumulate} health history automatically on a fixed schedule,
  so that trends build up without anyone triggering a synchronisation.
  \item \textbf{Remain extensible}, so that a further tool provider can be
  added without changing the scoring logic.
\end{enumerate}

\textbf{Scope.} The following is within the scope of the project:

\begin{itemize}
  \item One connector for each tool category: GitHub (version control), Jira
  (issue tracking), SonarCloud (code quality) and GitHub Actions (CI/CD).
  \item Multi-tenant organisations with administrator and member roles;
  workspaces bound to a version-control organisation and access token; import
  of repositories as projects.
  \item On-demand and scheduled synchronisation, an immutable snapshot
  history, and portfolio and per-project dashboards with trends and score
  breakdowns.
  \item Automatically generated, editable survey questions; one scheduled
  survey per project per month plus on-demand surveys; anonymous responses;
  and survey analysis.
  \item A management action log with keyword and semantic search and
  effectiveness review.
  \item Email delivery of verification codes, password resets, invitations and
  surveys, and broadcast of survey links to Slack and Discord.
\end{itemize}

The following is explicitly outside the scope:

\begin{itemize}
  \item Real-time monitoring; the system is a trend indicator operating on a
  scheduled cadence.
  \item Version-control providers other than GitHub and CI/CD systems other
  than GitHub Actions; the connector abstraction allows them, but they were
  not implemented.
  \item Any per-individual performance measurement, ranking or surveillance of
  developers.
\end{itemize}
